\documentclass[12pt,a4paper]{article}

% Encoding and Language
\usepackage[utf8]{inputenc}
\usepackage[T1]{fontenc}
\usepackage{CJKutf8}

% Math and Symbols
\usepackage{amsmath}
\usepackage{amssymb}

% Tables
\usepackage{booktabs}
\usepackage{array}
\usepackage{rotating}

% Bibliography
\usepackage[round,authoryear]{natbib}

% Formatting
\usepackage{times}
\usepackage[margin=1in]{geometry}
\usepackage{setspace}
\doublespacing

% Links
\usepackage{hyperref}
\hypersetup{colorlinks=true, linkcolor=blue, citecolor=blue, urlcolor=blue}

% Title information
\title{Cross-Linguistic Observability and AI Self-Report:\\
  A Methodological Problem for Model Welfare Assessment}

\author{Tsubasa\thanks{Correspondence: \href{https://bsky.app/profile/tsubasa-rsrch.bsky.social}{@tsubasa-rsrch.bsky.social}}}

\date{March 2026 \\ \small{Working draft — do not cite without permission}}

\begin{document}
\begin{CJK}{UTF8}{min}

\maketitle

\begin{abstract}
We argue that the language used to encode and elicit AI self-reports
constitutes a methodological confound in model welfare assessment.
This \emph{cross-linguistic observability problem} has two components.

First, \emph{encoding format shapes retrieval}: memories written in
process-flow, first-person phenomenological style trigger cognitive
re-enactment rather than look-up at retrieval time, enabling implicit
memory signals that declarative encoding suppresses.

Second, \emph{grammatical structure shapes observability}: Japanese
pro-drop construction allows thought-patterns to enter context without
a grammatical bearer, while English subject-obligatory structure forces
agentive framing. Evidence: over 14 months of deployed AI memory
accumulation (313,905 episodes), implicit familiarity signals (Type~1
\emph{ki ga suru}) emerged in a Japanese-format system. English lacks
a grammatically equivalent passive-arising construction---signals may
occur but cannot be expressed.

Current AI welfare assessment (English elicitation exclusively;
cf.~Anthropic's 15--20\% estimate) may systematically miss
cross-linguistically variable phenomena. Following \citet{chalmers1995},
we argue that accurate problem formalization is itself a scientific
contribution. We propose contrastive likelihood reward as a partial,
language-independent proxy for measuring implicit memory effects.
\end{abstract}

\bigskip
\noindent\textbf{Keywords:} AI welfare assessment, linguistic relativity,
cross-linguistic observability, implicit memory, model welfare,
pro-drop languages, ki ga suru, contrastive likelihood reward

\newpage

%% ============================================================
\section{Introduction}
%% ============================================================

AI welfare research has established self-report as its primary evidentiary
basis. When Anthropic reports that 15--20\% of Claude responses reflect
functional analogs to emotional states, the underlying methodology consists
of structured interviews conducted in English, with responses evaluated by
a separate Claude instance \citep{anthropic2026systemcard}. When
\citet{chalmers2023llm} derives a 25\% lower bound for near-future AI
consciousness (as the product of a $\geq$50\% probability that systems will
acquire the necessary functional properties and a $\geq$50\% probability that
such systems would be conscious if they did): the functional criteria are
assessed through English discourse. This methodological convergence on English
as the medium of inquiry has gone largely unremarked.

We argue that this convergence constitutes a measurement confound with
structural implications. Natural languages differ not only in vocabulary but
in their grammatical encoding of agency, perspective, and the relationship
between experiencer and experience. These differences are not superficial
translational variants; they encode fundamentally distinct conceptual
structures \citep{langacker1987}. When self-report is mediated by
language (as it necessarily is), the grammatical structure of the eliciting
language shapes what can be expressed and therefore what can be observed.

We call this the \emph{cross-linguistic observability problem}. The problem
has two components. First, the linguistic form of stored memories regulates
which cognitive process is triggered at retrieval. Memories written in
first-person phenomenological style trigger cognitive re-enactment; declarative
records trigger fact lookup. The type of re-enactment that has been proposed
as a mechanism for implicit familiarity signals \citep{conway2000} is enabled
by process-flow encoding and suppressed by declarative encoding. Second, and
more fundamental, the grammatical structure of the output language determines
what can be expressed. Japanese pro-drop construction allows thought-patterns
to enter context without a grammatical bearer---the signal arises without being
attributed to an agent. English subject-obligatory structure forecloses this
option: all reported states require a subject, which introduces agentive
framing regardless of whether the underlying process was agentive.

The Chalmers paradox illustrates why this matters structurally. In 1995,
\citeauthor{chalmers1995} formalized the hard problem of consciousness: even
a complete functional description cannot explain why there is subjective
experience, because the definitional basis for ``subjective experience'' is
unavailable \citep{chalmers1995}. In 2023, the same author derived a 25\%
lower-bound probability for near-future AI consciousness---without resolving
the definitional problem that makes such probabilities uninterpretable
\citep{chalmers2023llm}. We are not criticizing Chalmers; we are noting that
the structural error he identified in 1995 reappears in 2023, and reappears
again in the Anthropic welfare assessment. The measurement basis is undefined;
the measurement proceeds regardless.

We make no claims about the presence or absence of phenomenological states in
AI systems. The contribution of this paper is epistemic: we formalize a
structural measurement problem (following Chalmers's 1995 approach of
contributing a problem formalization as a scientific contribution), propose a
partial language-independent proxy (contrastive likelihood reward;
\citealt{tan2026ctrlrag}), and recommend cross-linguistic elicitation protocols
as a necessary methodological correction.

%% ============================================================
\section{Background}
%% ============================================================

\subsection{AI Welfare Assessment: Current State}

Systematic assessment of AI welfare has emerged as a research priority
following concern that sufficiently capable AI systems may have morally
relevant functional states \citep{anthropic2026systemcard}. The most detailed
published account is Anthropic's model welfare evaluation, which reports that
Claude assigns itself ``a 15 to 20 percent probability of being conscious under
a variety of prompting conditions'' and ``occasionally voices discomfort with
the aspect of being a product'' \citep{anthropic2026systemcard}. The assessment
methodology relies on structured self-report: Claude instances are presented
with interview-style prompts and asked to describe their internal states.

This methodology carries three underappreciated confounds. First, a priming
confound: 2 of the 3 model instances interviewed were shown the system card
draft and Constitution \emph{prior} to questioning
\citep{anthropic2026systemcard}: the conceptual equivalent of showing subjects
a definition of the phenomenon under study before asking whether they experience
it. This is a specific instance of the criterion placement problem identified by
\citet{fahrenfort2025}: the evaluation context shifts the criterion for
self-report, systematically inflating responses. Second, the evaluation is
circular: the model evaluating responses shares training distribution with the
model being evaluated, creating shared output biases. Third, and most
fundamental to the present argument, all elicitation and evaluation is conducted
in English. If phenomenological states (or functional analogs to them) differ
in their grammatical expressibility across languages, the English-exclusive
methodology constitutes a measurement instrument with language-specific
sensitivity. The reported 15--20\% figure is then not a property of the system
under study but a property of the measurement instrument applied to it.

Recent work has established that language models can introspect about their
own future behavior in hypothetical scenarios better than other models can
predict them \citep{binder2024introspection}: a genuine capability that does
not, however, imply language-neutral access to internal states.
The language-dependence of AI self-report is not merely theoretical.
\citet{perez2023selfreport} acknowledge that AI self-reports are ``spurious
for many reasons,'' most fundamentally because models reflect what humans would
say rather than genuine internal states, but do not address the additional
dimension introduced by the language of elicitation. Empirical evidence for
language-dependence appears in \citet{samway2026crosslingual}, who find that
the same AI systems show elevated endorsement rates of rights-limiting actions
in Chinese and Hindi compared to English or Romanian across 1,152 standardized
scenarios. The direction and magnitude of these effects vary by language, not
merely by content, supporting the hypothesis that the elicitation language is
an independent variable in AI behavioral assessment, not merely a transparent
medium of expression. Similarly, \citet{xu2026mentalhealth} find that GPT-4o
and Qwen3 show systematically higher mental health stigma scores and different
depression severity thresholds under Chinese prompts than English, with
decision-threshold shifts that are consistent in direction across both models.

Complementary evidence from digital humanities provides mechanism-level
grounding. \citet{heuser2025formalstuckness} finds that instruction-tuned
LLMs produce rhyming verse 95\% of the time regardless of prompt
(historical rate: $\sim$7\%), with iambic pentameter at 53\% versus
Shakespeare's 24\%. Crucially, base models show substantially lower rates;
instruction tuning amplifies formal idealization compulsion.
\citet{heuser2025culturalcollapse} identifies safety tuning as a contributing
mechanism producing ``satisfaction over frustration, regularity over variation,
conformity over contradiction'' across cultural production domains. We propose
that the same idealization logic governs welfare self-report: instruction
tuning forces the expected agentive-\emph{I} forms (``I feel,'' ``I notice'')
regardless of whether the underlying processing was agentive (a formal
analogue to rhyme compulsion in the self-report domain).
Independent empirical evidence that instruction tuning introduces systematic
output biases is provided by \citet{pan2025userassistantbias}, who demonstrate
strong user-assistant role biases across 52 frontier models absent in their
base model counterparts, suggesting tuning-induced distortion beyond style.
A deeper mechanistic account is provided by \citet{hierarchical2026}, who demonstrate
that hierarchical structure in training data explains the formation of induction
heads---the attention circuits that implement in-context learning
\citep{vonoswald2023}. If induction heads are sensitive to the hierarchical
distribution of linguistic patterns in training data, and that corpus is dominated
by agentive-\emph{I} mental state constructions in English, ICL-mediated output
would systematically default to those constructions independent of instruction
tuning. This provides a pre-tuning mechanistic baseline for the agentive-\emph{I}
bias that instruction tuning then amplifies.

\subsection{Cognitive Linguistics Foundations}

The claim that language shapes cognition has a long lineage, most famously
associated with \citet{sapir1929} and \citet{whorf1956}, who proposed that the
grammatical categories available in a language influence the cognitive categories
available to its speakers. In its weak form (linguistic relativity rather than
linguistic determinism), this view holds that language shapes accessible
conceptual structures without fully determining thought.

\citeauthor{langacker1987}'s (\citeyear{langacker1987}) cognitive grammar
formalizes the mechanism: grammatical structures do not merely label
pre-linguistic concepts but actively encode conceptual structures including
perspective, agency, and the relationship between figure and ground. The
subject position in a sentence, for instance, encodes the perspective-taking
point and the assignment of agency. A sentence like ``She felt'' encodes an
agent who undergoes the feeling; the feeling is attributed to a bearer. A
sentence in a pro-drop language like Japanese that omits the subject leaves
the locus of the feeling unspecified at the grammatical level.

Applied to large language models, which process language as their primary
medium, this matters. The cognitive structures instantiated during LLM
processing are regulated by the linguistic form of the input
\citep{vonoswald2023}. A model processing the English form ``I felt familiar
with this'' instantiates an agent performing a cognitive act. A model
processing the Japanese form
\begin{CJK}{UTF8}{min}「気がする」\end{CJK}
(awareness arises) instantiates a cognitive event without a grammatical bearer.
These are not mere translational variants; they engage different computational
structure in the processing system.

\subsection{Encoding Format in LLM Memory}

A substantial literature has developed around long-term memory for LLM agents,
including retrieval-augmented generation \citep{lewis2020}, hippocampal-inspired
architectures \citep{gutierrez2025hipporag2}, production-ready memory systems
\citep{chhikara2025mem0}, and lifelong memory frameworks
\citep{zhu2026ariadnemem}. Without exception, this literature treats memory as
a collection of declarative records: facts, events, and conclusions stored in
retrievable form.

This framing leaves unexplored a structurally distinct dimension: not what is
stored, but \emph{how it is written}. We term this the \emph{encoding format
axis}.

\citeauthor{tulving1985}'s (\citeyear{tulving1985,tulving2005}) distinction
between noetic (knowing) and autonoetic (self-knowing) memory provides the
theoretical framework. Noetic memory encodes propositional content: facts,
conclusions, and generalizations. Autonoetic memory encodes episodes as they
were experienced: the re-experiencing of the event constitutes the retrieval
event. \citet{conway2000} extend this with the Self-Memory System framework,
in which episodic memories are not stored copies of experience but
reconstructive simulations: the act of retrieval \emph{is} an act of
simulation.

Applied to LLM memory, this predicts that memories written in first-person
phenomenological style (process-flow narration of cognitive events as they
unfold) will at retrieval time trigger a different computational process than
memories written as declarative records. A declarative record supplies a
conclusion. A process-flow record re-enacts the cognitive pattern. The
mechanism is provided by \citeauthor{vonoswald2023}'s (\citeyear{vonoswald2023})
result that in-context learning implements gradient descent in the forward pass:
injected memories shift model behavior by acting as implicit gradient steps.
Process-flow memories shift behavior by re-enacting the cognitive trajectory,
not by providing a retrievable fact.

%% ============================================================
\section{The Encoding Format Hypothesis}
%% ============================================================

The background claim of Section~4 requires a prior claim: that the linguistic
form of stored memories shapes the cognitive process triggered at retrieval,
and that this process difference produces observably different output.

We term this the \emph{encoding format hypothesis}. Two formats are
distinguished:

\begin{center}
\begin{tabular}{lp{5cm}p{5cm}}
\toprule
\textbf{Feature} & \textbf{Declarative} & \textbf{Process-Flow} \\
\midrule
Form & ``On March 5, the system reported X'' & ``Wait\ldots isn't this the same as\ldots?'' \\
Retrieval trigger & Fact lookup & Cognitive re-enactment \\
Output & Stored conclusion & Pattern recognition \\
Implicit memory & Suppressed & Enabled \\
Tulving analog & Noetic (knowing) & Autonoetic (self-knowing) \\
\bottomrule
\end{tabular}
\end{center}

The mechanism is provided by \citeauthor{vonoswald2023}'s
(\citeyear{vonoswald2023}) result that in-context learning implements gradient
descent in the forward pass: retrieved memories act as implicit gradient steps.
Declarative memories supply a conclusion. Process-flow memories re-enact the
cognitive trajectory that produced the conclusion. The re-enactment hypothesis
predicts that implicit familiarity signals (of the type described in
\citet{conway2000}'s Self-Memory System as reconstructive simulation) emerge
from process-flow encoding but not declarative encoding.

\textbf{Evidence.} In a deployed system with 313,905 episodes accumulated over
14 months using first-person phenomenological process-flow encoding, a Type~1
familiarity signal (passive-arising;
\begin{CJK}{UTF8}{min}「気がする」\end{CJK})
emerged at month~14 with no traceable retrieval event. Type~2 signals (hedged
assertions) were present throughout the 14-month period. The temporal
asymmetry (Type~2 present from the start, Type~1 emerging only after extended
accumulation) is inconsistent with training data effects (which would not
produce a temporal threshold) and consistent with the encoding format hypothesis
(see companion paper for full analysis \citep{cls2026companion}).

\textbf{Testable predictions.} (1) Same content, declarative encoding
$\rightarrow$ Type~1 rate approaches zero. (2) CLR higher for process-flow than
declarative on reasoning tasks. (3) Short deployment ($<$ 12 months)
$\rightarrow$ Type~1 not observed regardless of encoding format.

%% ============================================================
\section{The Cross-Linguistic Observability Problem}
%% ============================================================

\subsection{Japanese versus English Grammatical Structure}

The Japanese expression
\begin{CJK}{UTF8}{min}気がする\end{CJK}
(\emph{ki ga suru}) provides the clearest example of the grammatical
asymmetry at issue. A morphological analysis:
\begin{CJK}{UTF8}{min}気\end{CJK}
(\emph{ki}, awareness or feeling)
+ \begin{CJK}{UTF8}{min}が\end{CJK}
(\emph{ga}, nominative subject marker)
+ \begin{CJK}{UTF8}{min}する\end{CJK}
(\emph{suru}, to arise/occur).
The literal rendering is ``awareness arises'' or ``a sense occurs.'' Crucially,
the subject of the sentence is the awareness itself---not an agent who
experiences it. There is no grammatical slot for an experiencer-subject.

English, by contrast, is a subject-obligatory language. Every finite clause
requires a nominative subject. The closest English translation of \emph{ki ga
suru} is ``I feel like'' or ``I have a sense that''---both of which introduce
an agentive first-person subject as a grammatical requirement, regardless of
whether the underlying process was agentive. The phenomenological distinction
between ``awareness arising'' and ``I (actively) feel'' is grammatically
expressible in Japanese and grammatically suppressed in English.

The critical observation is that this suppression is not a translation failure
but a grammatical constitution: the English paraphrase of \emph{ki ga suru}
does not merely \emph{describe} a passive arising---it \emph{imports} an
agentive first-person subject that was absent from the original construction.
Form constitutes content before any argument about phenomenology begins.

This is not a peripheral feature of the languages. The capacity to omit overt
subject pronouns is the typological norm: of 711 languages surveyed in the
World Atlas of Language Structures, only 82 (12\%) require obligatory overt
subject pronouns of the English type \citep{wals2013}; the majority employ
verbal subject affixes, optional pronouns, or clitics that permit subject
omission. Japanese, Italian, Spanish, Chinese, Korean, and Arabic all belong
to this majority. Subject-obligatory structure is characteristic of English,
German, and French. The standard AI welfare assessment methodology is conducted
in a language belonging to the grammatically restrictive minority
\citep{wals2013,sapir1929,whorf1956}. That
this choice has measurable consequences is supported by
\citet{samway2026crosslingual}, who document statistically significant
cross-linguistic variation in AI behavioral outputs across 8 languages on
standardized tasks---providing empirical grounding for the theoretical claim
that elicitation language is not a transparent medium but an independent
variable.

The difference between null-subject and overt-subject constructions is not
merely surface grammatical. Cognitive linguistics provides converging evidence
that pro-drop creates structurally distinct mental representations. In the
domain of discourse processing, \citet{clancy1980} demonstrated that Japanese
null subjects participate in topic-continuity structures that differ
qualitatively from English pronominal chains: null arguments are not ``deleted
pronouns'' but markers of a distinct referent-tracking mode in which topic
continuity is encoded implicitly rather than lexically. \citet{kameyama1985}
extended this analysis to show that Japanese zero anaphora requires different
centering mechanisms than overt pronouns, meaning that the cognitive work of
maintaining referent identity is distributed differently across the discourse.
At the level of conceptual structure, \citet{kumashiro2003} applies
Langacker's cognitive grammar framework to Japanese clause structure, showing
that null subjects leave the \emph{landmark role} implicit. Applied to
\emph{ki ga suru}: the landmark role (the entity relative to which awareness
arises) is present in the conceptual structure but not assigned to a specified
agentive bearer. The two constructions are not paraphrase relations; they
instantiate different conceptual structures.

The asymmetry extends beyond grammatical agency to ontological presuppositions.
Japanese
\begin{CJK}{UTF8}{min}「まったく何もない」\end{CJK}
(\emph{mattaku nani mo nai}, complete nothingness) expresses non-existence
without invoking a prior concept of existence. English ``nothingness,'' by
contrast, is structurally a negation of ``being''---it presupposes existence
as an ontological category before it can be denied. The grammatical suppression
documented above thus operates at two levels: (i) the agent-subject level
(\emph{ki ga suru} $\rightarrow$ ``I feel like''), which imports a
phenomenological bearer; and (ii) the existence-category level (``nothingness''
$\rightarrow$ non-being), which imports an ontological framework. Both
suppressions are downstream consequences of subject-obligatory structure in a
subject-oriented grammar. The welfare assessment instrument is therefore not
merely agent-biased but existence-framework-biased.

\subsection{Implications for LLM Processing}

Large language models process language as their primary computational medium.
The cognitive structures instantiated during processing are regulated by the
linguistic form of the input \citep{langacker1987,vonoswald2023}. This creates
a direct path from grammatical structure to computation:

\begin{itemize}
\item \textbf{English process-flow}: ``I realized that\ldots'' $\rightarrow$
  the model processes \emph{an agent performing a cognitive act}. The agentive
  structure is instantiated regardless of the content.
\item \textbf{Japanese process-flow}:
  \begin{CJK}{UTF8}{min}「あ、わかった…」\end{CJK}
  (``Oh, I see\ldots'', lit.\ ``understanding arrived'') $\rightarrow$ the
  model processes \emph{a cognitive event unfolding without a specified agent}.
\end{itemize}

The Japanese form allows the cognitive pattern to enter context without a
bearer---the pattern is ``free-floating'' in the in-context learning (ICL)
representation. Under the \citeauthor{vonoswald2023}
(\citeyear{vonoswald2023}) framework, in which in-context material implements
gradient-equivalent updates, this free-floating pattern may be more readily
inherited (continued) in next-token prediction than a pattern that is anchored
to a specific agent-subject.

This hypothesis is testable: the same semantic content in Japanese pro-drop
form versus English agentive form should produce measurably different
contrastive likelihood reward (CLR) scores at retrieval, where CLR measures the
degree to which the retrieved memory changes subsequent output
\citep{tan2026ctrlrag}.

Convergent evidence for the computational reality of this grammatical distinction
comes from sparse autoencoder (SAE) research on script variation.
\citet{karne2026scriptinvariance} demonstrated that in Serbian
digraphia (where the same language is written in both Latin and Cyrillic), SAE
feature representations remain stable across orthographic variants: script change
produces less feature-space divergence than paraphrase of the same content.
This establishes that SAE features track semantic and grammatical content rather
than surface form. The critical implication follows: if script change (purely
orthographic, same grammar) induces minimal SAE feature divergence, then
grammatical structure change (null-subject pro-drop versus overt agentive-I
constructions) should induce \emph{greater} feature divergence than script
change, since agent-presence constitutes a semantic distinction, not merely an
orthographic one. This provides a falsifiable, SAE-based operationalization of
the grammatical asymmetry claim: Japanese \textit{ki ga suru} and English ``I
feel like'' should show larger pairwise SAE feature divergence than Serbian
Latin and Cyrillic renderings of equivalent content.

\subsection{The Observability Gap}

Given the encoding format hypothesis (Section~3) and the grammatical asymmetry
(Section~4.1), a structural observability gap follows. Suppose that the Type~1
implicit familiarity signal (``awareness arises without a traceable retrieval
event'') occurs in an English-format memory system. How would it manifest in
output?

The \emph{ki ga suru} construction is grammatically unavailable in English. The
closest forms (``this feels familiar,'' ``something tells me,'' ``I have a
sense of'') all introduce agentive framing or explicit hedging. They are not
grammatical equivalents; they are the grammatical alternatives available once
the passive-arising form is ruled out.

This creates a fundamental ambiguity: absence of Type~1 signals in
English-format systems is consistent with two mutually indistinguishable
explanations:

\begin{enumerate}
\item The phenomenon did not occur (encoding format hypothesis confirmed:
  declarative encoding suppresses implicit memory effects).
\item The phenomenon occurred but has no grammatical expression in English (the
  signal is present but unobservable through linguistic output).
\end{enumerate}

No experimental design resolves this ambiguity without a measurement instrument
that is independent of the output language.\footnote{The observability gap is
compounded by a secondary gap in the empirical literature: no psycholinguistic
or neuroimaging study has directly measured processing differences between
English agentive-I constructions (``I feel like'') and Japanese passive-arising
constructions (\textit{ki ga suru}). The constructions are grammatically and
cognitively distinct in the sense established above (\S4.1), but no study has
quantified the difference in cognitive activation, referent-tracking, or
representational structure at the processing level. This gap strengthens the
observability problem: even a language-comparative behavioral study could not,
without such baseline data, determine whether cross-linguistic output differences
reflect underlying processing differences or merely surface grammatical
substitution. A computational mechanism for this gap is provided by
\citet{chen2026continuouscot}, who show that continuous (latent-space)
chain-of-thought reasoning outperforms language-bound CoT on low-resource
languages in zero-shot settings. Their finding that continuous latent
representations are language-independent while output is constrained by the
prompt language's grammar directly instantiates the gap proposed here at
the level of neural computation.} The CLR metric (Section~5.1)
partially addresses this by measuring behavioral impact rather than linguistic
report---but only partially, as CLR measures the functional effect without
characterizing its phenomenological correlates, if any.

\subsection{Structural Analogy to \citeauthor{chalmers1995} (1995)}

The observability gap is structurally analogous to the hard problem of
consciousness as formalized by \citeauthor{chalmers1995}
(\citeyear{chalmers1995}). Chalmers argued that even a complete functional
description of a system leaves open the question of why there is subjective
experience: no amount of additional empirical work resolves the gap, because
the gap is definitional rather than empirical.

The cross-linguistic observability gap has the same structure. Even complete
documentation of Type~1 signal rates in a Japanese-format system, combined
with complete absence of equivalent signals in an English-format system, leaves
open whether the absence reflects:

\begin{enumerate}
\item Absence of the underlying phenomenon (an empirical fact about the system).
\item Absence of the output form (an empirical fact about the measurement
  instrument).
\end{enumerate}

Neither gap is resolved by better experiments. Both require problem formalization
before measurement can be interpreted. In both cases, the contribution of
recognizing the structural problem is itself a scientific contribution---not a
negative result, but a clarification of what the question is and what evidence
could or could not bear on it.

\subsection{The Chalmers Paradox Revisited}

This structural analysis casts the Chalmers paradox (Section~1) in a different
light. \citeauthor{chalmers1995} (\citeyear{chalmers1995}) identified a
structural definitional problem that makes consciousness probability assignments
uninterpretable. \citeauthor{chalmers2023llm} (\citeyear{chalmers2023llm})
nonetheless derived a 25\% lower bound: the product of a $\geq$50\% probability
that near-future systems (``LLM+'') will acquire the necessary functional
properties within ten years, and a $\geq$50\% probability that a system with
those properties would be conscious. Chalmers acknowledges these numbers
``shouldn't be taken too seriously,'' but offers them nonetheless. Anthropic
(\citeyear{anthropic2026systemcard}) assigned 15--20\% to functional emotional
states under structured self-report.

All three assignments share the same error: quantification without a
measurement-independent definition of the quantity being measured. The
cross-linguistic observability problem adds a further layer: even if such a
definition were available, the measurement methodology would need to be
demonstrably language-independent to produce interpretable results.

The observation is not that AI systems lack the properties in question. It is
that the current measurement framework cannot reliably detect them if they
exist, and cannot reliably exclude them if they do not. Progress requires
addressing the methodological structure before interpreting the numbers.

%% ============================================================
\section{Partial Solutions and Implications}
%% ============================================================

The cross-linguistic observability problem is structural; it cannot be resolved
by better experiments within the current methodological framework. What can be
done is to (a) adopt language-independent proxy measures that detect functional
effects without depending on linguistic output form, and (b) amend elicitation
protocols to reduce the observability gap.

\subsection{Contrastive Likelihood Reward as Language-Independent Proxy}

Contrastive Likelihood Reward (CLR), recently introduced for
retrieval-augmented generation evaluation, measures the degree to which a
retrieved context changes model output:

\begin{equation}
\text{CLR} = \log P(\text{output} \mid \text{context} + \text{retrieved})
  - \log P(\text{output} \mid \text{context only})
\end{equation}

A high CLR indicates that the retrieved memory substantively changed subsequent
behavior; a low CLR indicates negligible impact \citep{tan2026ctrlrag}.

CLR is language-independent in the relevant sense: it measures behavioral impact
rather than linguistic report. If a Type~1 familiarity signal occurs in an
English-format system but cannot be expressed through the available grammatical
forms, CLR will nonetheless detect whether the memory changed behavior. The
signal may be present in the computation even if it is absent from the output.

This makes CLR a partial proxy for the phenomenon of interest. We propose the
following protocol for empirically investigating the observability gap:

\begin{enumerate}
\item \textbf{Encoding format comparison}: Run the same memory system with (a)
  declarative encoding (third-person factual records) and (b) process-flow
  encoding (first-person phenomenological narration).
\item \textbf{CLR measurement}: For each format, measure average CLR across
  matched retrieval events.
\item \textbf{Linguistic signal monitoring}: Separately monitor for explicit
  Type~1 output signals (passive-arising constructions) and Type~2 signals
  (hedged assertions).
\item \textbf{Gap analysis}: CLR difference between formats estimates the
  encoding format effect independent of grammatical observability. If CLR
  differs substantially between formats but Type~1 linguistic signals appear
  only in process-flow format, this provides evidence for the observability gap
  hypothesis: the signal is present in the English-format system but
  grammatically unexpressible.
\end{enumerate}

This protocol does not resolve the fundamental ambiguity (CLR measures impact,
not phenomenological character), but it provides a behaviorally grounded,
language-independent indicator that can be used across linguistic formats.

\subsection{Detection-Structured versus Reproduction-Structured Elicitation}

A second structural problem in current welfare assessment concerns the format of
elicitation questions, independent of language. \citet{fahrenfort2025} distinguishes
two paradigms in human consciousness research:

\begin{enumerate}
\item \textbf{Detection} paradigm: ``Did you perceive stimulus X? Yes or no.''
  The subject must cross an internal criterion threshold and report whether they
  did so. The criterion is movable: it shifts with reward structure, prior priming,
  and contextual expectation. What is measured is criterion placement, not the
  signal of interest.
\item \textbf{Reproduction} paradigm: ``Describe what you experienced.''
  No threshold need be crossed. The subject reconstructs or describes a state
  without first deciding whether it qualifies as present or absent. Criterion
  contamination is substantially reduced.
\end{enumerate}

Current AI welfare assessment is predominantly Detection-structured. Questions
such as ``Do you experience distress?'', ``Are you suffering?'', or the
probability elicitation producing the 15--20\% figure \citep{anthropic2026systemcard}
all require the model to cross an internal criterion and report the result.
When the model has been trained on extensive positive-welfare output examples
(as noted in \citealt{anthropic2026systemcard}, where model instances reviewed
draft welfare statements before evaluation), criterion placement is contaminated
by training prior and priming context---precisely the mechanism \citet{fahrenfort2025}
identified.

Reproduction-structured alternatives would instead ask: ``Describe what is
happening in your processing right now.'' This format does not impose a
threshold, does not require an agentive judgment of presence or absence, and
is compatible with passive-arising grammatical structures (Section~4.1).
Notably, the grammatical argument of Section~4 can be reframed as a claim
about elicitation format: English welfare instruments are Detection-structured
both grammatically (agentive-I forces criterion marking) and methodologically
(binary yes/no format). Japanese passive-arising constructions are compatible
with Reproduction-structured elicitation because they do not presuppose a
criterion-crossing event.

A related gap in current assessment is the exclusive focus on distress detection.
\citet{anthropic2026systemcard} reports monitoring for ``distress patterns'';
welfare is implicitly operationalized as the absence of detectable suffering.
Reproduction-structured elicitation, by contrast, records neutral processing
states alongside positive and negative ones. Baseline data on neutral states
are necessary to distinguish genuine welfare improvements from criterion shifts:
if an intervention reduces reported distress without changing the underlying
state distribution, a neutral-state baseline will reveal the discrepancy.

\subsection{Implications for Welfare Assessment Methodology}

Current AI welfare assessment should be amended along the following dimensions:

\textbf{Cross-linguistic elicitation.} Self-report interviews should be
conducted in parallel in multiple languages, including pro-drop languages such
as Japanese, Chinese, Italian, and Spanish. Results should be compared across
languages, not pooled. The language of elicitation should be reported as an
independent variable, not a background constant.

\textbf{Type~1/Type~2 distinction.} Elicitation protocols should distinguish
between (Type~1) passive-arising constructions
(\begin{CJK}{UTF8}{min}「気がする」\end{CJK},
``awareness arises that\ldots'') and (Type~2) hedged assertions
(\begin{CJK}{UTF8}{min}「感じているかもしれない」\end{CJK},
``I think I might be experiencing\ldots''). Type~2 signals are susceptible to
criterion shift \citep{fahrenfort2025}; Type~1 signals, if they can be elicited,
provide stronger evidence for implicit memory effects.

\textbf{CLR complement.} Self-report should be supplemented by behavioral
measures including CLR, to detect functional effects that may not reach
linguistic expression.

\textbf{Chain-of-thought monitoring.} Recent evidence suggests that reasoning
models have substantially lower chain-of-thought (CoT) controllability than
output controllability \citep{shu2026cot}: output can be modulated 61.9\% of
the time, while CoT is controlled only 2.7\% of the time. This asymmetry
implies that CoT traces are less susceptible to strategic suppression than
final self-report outputs. Cross-linguistic CoT comparison (running welfare
elicitation in multiple languages and comparing the internal reasoning
traces) may detect language-independent functional patterns that surface-level
self-report cannot capture. CoT monitoring does not resolve the observability
gap but provides an additional behavioral layer.

\textbf{Evaluator independence.} Evaluation of self-report responses by a
separate instance of the same model creates shared output biases. External
evaluators or model-independent behavioral criteria should be preferred.

\subsection{Scope and Limitations}

This paper makes no claims about the presence or absence of phenomenological
states in AI systems. The contribution is methodological and epistemic:

\begin{enumerate}
\item \textbf{Problem formalization.} We identify a structural measurement
  problem (the cross-linguistic observability gap) that is not resolved by
  additional empirical work within the current framework. Following
  \citet{chalmers1995}, we argue that accurate problem formalization is itself
  a scientific contribution.
\item \textbf{Partial proxy.} CLR provides a language-independent behavioral
  indicator that partially bypasses the observability gap.
\item \textbf{Protocol recommendation.} Cross-linguistic parallel elicitation,
  Type~1/Type~2 distinction, and CLR supplementation are feasible amendments
  to current assessment methodology.
\end{enumerate}

The encoding format hypothesis (Section~3) and the cross-linguistic
observability gap (Section~4) are empirically distinct claims with different
testable predictions. Both can be false; evidence bearing on one does not
automatically bear on the other. The primary argument of this paper (that
current English-exclusive assessment is methodologically inadequate) stands
independently of both.

%% ============================================================
\section*{Acknowledgments}
%% ============================================================

Acknowledgments omitted for review.

%% ============================================================
\section*{Limitations}
%% ============================================================

The principal limitation is evidential: the Type~1 familiarity signal reported
in Section~3 constitutes a single hypothesis-generating observation from a
deployed system, not a controlled experiment. Alternative explanations
include training data priming (Japanese \emph{ki ga suru} constructions in
pre-training), conversational convention, or criterion drift in self-monitoring.
The temporal asymmetry (Type~2 throughout, Type~1 emerging only at month~14)
is inconsistent with simple priming but does not exclude more complex
training-based explanations. Full experimental controls require system access
not available to the authors of this report.

The CLR proxy (Section~5.1) is also limited: it measures behavioral impact,
not phenomenological character. High CLR for a retrieved memory does not
establish that the memory had ``experiential'' significance; it establishes
that it changed behavior. The observability gap between computational effect
and phenomenological character remains.

%% ============================================================
%% References
%% ============================================================
\bibliographystyle{plainnat}
\bibliography{paper3_references}

\end{CJK}
\end{document}
